\uuid{NedA}
\titre{Calcul de dérivée}
\auteur{bodin ; hanani ; mzari}
\organisation{exo7}
\licence{CC BY-NC-SA 4.0}
\datecreate{2020-01-20}
\niveau{L1}
\module{Analyse}
\chapitre{Dérivabilité des fonctions réelles}
\sousChapitre{Autre}
\difficulte{1}
\theme{Dérivées, calcul de dérivée}
\exo7qcmid{lille-1.derivees.006}

\contenu{
\begin{question}
\qclassification{124.00}
\qtags{motcle=calcul de dérivée}
Quelles sont les bonnes réponses ?
\begin{answers}
    \bad{La dérivée de $f(x)=\sin [(2x+1)^2]$ est $f'(x)=2\cos [(2x+1)^2]$.}
    \good{La dérivée de $f(x)=\sin [(2x+1)^2]$ est $f'(x)=4(2x+1)\cos [(2x+1)^2]$.}
    \good{La dérivée de $f(x)=\tan (1+x^2)$ est $\displaystyle f'(x)=\frac{2x}{\cos ^2(1+x^2)}$.}
    \bad{La dérivée de $f(x)=\tan (1+x^2)$ est $\displaystyle f'(x)=1+\tan ^2(1+x^2)$.}
\end{answers}
\begin{explanations}
De manière plus générale, $(\sin u)'=u'\cos u$ et
$$(\tan v)'=\frac{v'}{\cos ^2v}=v'(1+\tan ^2v).$$
Il suffit de prendre $u=(2x+1)^2$, $v=1+x^2\Rightarrow u'=4(2x+1)$ et $v'=2x$.
\end{explanations}
\end{question}
}
